% abstract.tex — unstructured summary of the model and the four findings.

\begin{abstract}
\noindent
When a sensory cue signals that some locations are more valuable than
others, an ideal observer can adapt in more than one way: by shifting its
decision criteria, by re-allocating spatial attention to sharpen
perception where it matters, or---when sensory noise is correlated across
locations---by decorrelating the encoded population code. We develop a
normative model of value-directed attention in a cued change-detection
task that places these three levers within a single optimisation, with
the cross-location correlation a free parameter whose zero limit recovers
the independent-noise case. Four results follow. Criterion adjustment is
typically the dominant route to encoding value, capturing a median of
roughly three-quarters of the available reward gain across the parameter
space, with a substantial tail in which attention contributes materially.
The benefit of attention is non-monotonic in the benefit/cost ratio of
enhancement to suppression, peaking in an intermediate band whose lower
edge we give in closed form, and it is concentrated in a graded
regime---low cue validity, high value contrast, and cost-dominant
asymmetry---which we map quantitatively and turn into validity thresholds
for cueing designs. Finally, optimal attention never turns away from the
high-value location under a predictive cue, a result that holds as a
conditional theorem; under a counter-predictive cue, inverted allocation
becomes optimal, a new and falsifiable prediction.
\end{abstract}
